Fix the admissible precision. Use the witness constructed in the
proof of simultaneous family selection with the staged bounds \(K_*\).
Choose the cloud exponent \(c_D\) at least \(805\), then the node
exponent \(c_h\) at least \(1000\), in each position taking the maximum
with the other bounds there. The family \(\mathfrak F\) is constructed
with these constants. Its selected period and trace records give
the two section systems used below.
\releasecite{Development/AbelianVarietyTheory/Proofs/GlobalInterpolation/SzRowDelTOGap.lean}{12}{Cloud exponent}
\releasecite{Development/AbelianVarietyTheory/Proofs/GlobalInterpolation/SzRowNodeGapBlockLength.lean}{61}{node exponent}
\releasecite{Development/AutomorphicArithmetic/CharacterBoxes/DesignS.lean}{29}{order of choices}.

The simultaneous choice gives one class and one bundle selector for
the period, twist, divisor and row conditions. A built point carries
its own numerical record and its weight certificate. Membership in
the selected class identifies the triple of that record, its auxiliary
object and its bundle by a dependent equality with the selected triple.
Thus the weight certificate and the properties of the bundle transport
along the same equality. The final change of division data carries
both sample arrays and the weight certificate. Geometric compatibility
supplies comparison data on each bundle, giving one compatible
selector. The resulting point retains the bundle's period face,
summands and columns. The primitivity, divisor, goodness and class-sum
inputs give the point conditions. The geometric comparison gives the
additional condition called \texttt{More} on that point. These two
conditions are the inputs for the construction of the two runs at the
selected index.
\releasecite{Development/AbelianVarietyTheory/Definitions/GlobalTwisting/ValidForConstructedPointClassesEventuallyWorking.lean}{250}{Built triple and selection}
\releasecite{Development/AbelianVarietyTheory/Definitions/GlobalTwisting/SetUnipotentTranslate.lean}{381}{transport of the certificate}
\releasecite{Development/AutomorphicArithmetic/ThetaCorrespondence/ThetaPairLifts/ThirdPlaceSeries/NonemptyHypsOfAllRows.lean}{49}{joint class and selector}
\releasecite{Development/AutomorphicArithmetic/ThetaCorrespondence/ThetaPairLifts/ThirdPlaceSeries/PartSel.lean}{32}{data retained by the comparison}
\releasecite{Development/AutomorphicArithmetic/ThetaCorrespondence/ThetaPairLifts/ThirdPlaceSeries/TupAllUsedAtPart.lean}{46}{point conditions and their use}.

Take one threshold \(x_*\) dominating the finitely many thresholds of
the retained comparisons and the family's covering threshold. At any
real scale \(x\in\operatorname{Win}(\mathfrak F)\) above \(x_*\),
choose a seed \(i_0\in I_x\). In the classical window argument,
use the realized seed with the admissible pair of sample arrays.
For each \(u\in H\) the covering
property gives an index \(i(u)\in I_x\) with translate \(u\) and the
seed's \(a_t,b_t\) under the common numerical labels. The common word
labels preserve their bits. The constructions below apply to the seed
and to all these representatives, using each index's own periods,
tests and columns. The zero representative \(i(0)\) may differ from
\(i_0\). The windows contain arbitrarily distant intervals \([Y,Y^{3/2}]\)
above this threshold.
\releasecite{Development/AutomorphicArithmetic/CharacterBoxes/DesignS.lean}{55}{Intervals before the scale is chosen}
\releasecite{Formalism/SharpRestriction/FamCore.lean}{110}{word bits}
\releasecite{Formalism/SharpRestriction/FamCore.lean}{132}{all seed-matched translates}
\releasecite{Development/AutomorphicArithmetic/ThetaCorrespondence/ThetaPairLifts/ThirdPlaceSeries/RowRRowsOfParts.lean}{26}{uniform row thresholds}.

At each such index the numerical certificate gives, with
\(\Delta=\lceil x^{c_D}\rceil\) and \(H_{\rm node}=2\lceil x^{c_h}\rceil\),
\[
  v_a=4\Delta(k-2-a),\qquad
  b_{\tau,l}=B_\tau+lH_{\rm node},\qquad
  e_{\tau,l}=b_{\tau,l}+\tfrac12 H_{\rm node}
  \quad(0\le a\le k-2,\ 0\le l<m).
\]
For \(0\le a\le k-3\), the interior entry in position \(a+1\) of a
word is \(v_a-\tfrac12-\epsilon_a\Delta\), where \(\epsilon_a\) is its bit.
The consecutive cloud gap is \(4\Delta\), the bit displacement is
\(\Delta\), and the intervening node has displacement \(H_{\rm node}/2\).
The top start is \(b_t+\lfloor b'/x\rfloor\). A thick start belongs to
\([\lceil b\rceil,\lceil b\rceil+|\mu(E)|-1]\) with the prescribed
torsion condition, and a thin start is the prescribed \(\lceil w''\rceil\).
The node allowance is \(\eta_{\rm wt}=mH_{\rm node}\).
All words use these two spacings in the
\hyperlink{proof.section-run-design}{section weight estimates}.
At stage \(j=1,\ldots,m\) the signature blocks have sizes \(k+j-1\)
and \(m-j\). At each embedding the first tableau reads the entries
\(a_+^Q\) of the selected record, and the second reads \(-a_-^Q\)
in reverse order. These are the same entries as in the geometric
section weight, with common copy number \(r+t\). Thus the row weight
equals \(\kappa_j\) at every stage, with an empty second block at
\(j=m\).
\releasecite{Development/AbelianVarietyTheory/Definitions/GlobalTwisting/PairRowIffQdRows.lean}{93}{Exact weight certificate}
\releasecite{Formalism/LocalRingCardinality/RhoDetPrimeEqMuTildeMulChiPow.lean}{91}{word arrays}
\releasecite{Formalism/LocalRingCardinality/DataKap.lean}{20}{entries of both tableaux}
\releasecite{Development/AbelianVarietyTheory/Proofs/DivisorProofs/ExistsIsRec.lean}{24}{equality of tableaux and copy numbers}
\releasecite{Development/AutomorphicArithmetic/ThetaCorrespondence/ThetaPairLifts/ThirdPlaceSeries/RunModelArch.lean}{55}{ordinary and terminal signatures}.

Now work at one selected index. Write \(Q=g_0g_1\). Its marked
polarized core gives a common base level \(Q\), a single level action
and two covers at \(Qp_o\), with distinct own primes
\(p_o=p_{\rm de},p'_{\rm de}\). The refinement at \(p_o\) uses its
marked half and fixed determinant character \(\delta_{p_o}\). At a
finite place \(v\) use the run with \(p_o(v)=p'_{\rm de}\) if
\(v\mid p_{\rm de}\), and with \(p_o(v)=p_{\rm de}\) elsewhere. At
either refinement the coefficient
field equals the base field. Changing the copy number changes the
coefficient function on the same cover.
\releasecite{Development/AutomorphicArithmetic/ThetaCorrespondence/ThetaPairLifts/ThirdPlaceSeries/ExistsRunsAt.lean}{18}{Both runs with their joint action}
\releasecite{Development/AbelianVarietyTheory/Proofs/HeightProofs/ExistsSpecialPrimeLevelCoverOfFieldConditions.lean}{41}{field of the refinement}
\releasecite{Development/AbelianVarietyTheory/Definitions/Jets/EmbShortThLeThetaPlain.lean}{72}{the cover at other copy numbers}.

Let \(K_s\) be the embedded geometric field called \(K_0\) in the
section estimate. The preliminary column witness pins its exact
value. On the assembled stages the point's field comparison gives
\[
 K_s=\operatorname{FO}(D)=\operatorname{image}_{\mathbb C}(\text{\texttt{Base.Ks}})
    =\operatorname{listField}(\text{field data},\text{fixed frame}).
\]
The equality of embedded fields follows by identifying the images
of the rings of integers and taking fractions. The list includes
the embedded \(E\), the marked CM field \(M'\), the ray field of
\(E\), the specified roots of unity, the unramified finite readings
of \(\mu_{0,j},\eta,\delta,\delta'\), the finite \(\omega\)-values,
and the type-norm fields of all stages. For \(\eta\) the reading at
\(a\) is \(\eta(a)(\mathrm{N}a)^{-1/2}\). The two finite-order
characters \(\delta,\delta'\) also have their required idele values
in this field. The ray modulus \(Qp_{\rm de}p'_{\rm de}\) contains
both \(Qp_o\) as divisors, giving the required Weber values. The fixed
character readings thus occur before adjoining the moving characters.
The marking field \(K_{\rm mark}\) (\texttt{LL0}) contains the values
of the separate character \(\chi_\Psi\). Both embedded fields persist
in the enlargement to the period field \(K\) (\texttt{LL}), with
\(K_s\subseteq K_{\rm mark}\subseteq K\). The identity for \(K_s\)
is at the given index and allows the fields to vary with the translates.
\releasecite{Development/AbelianVarietyTheory/Proofs/HeightProofs/CoreFrameCompatibleConstructedPoint.lean}{11}{Exact field pin at the constructed point}
\releasecite{Development/AbelianVarietyTheory/Proofs/IwasawaDecomposition/IndexBoundCoefficientFieldOfExactSmallField.lean}{28}{assembled field and both level pins}
\releasecite{Development/AbelianVarietyTheory/Proofs/HeightProofs/FieldSmallIffOpo.lean}{10}{embedded integral readings and fractions}
\releasecite{Development/AbelianVarietyTheory/Definitions/AutomorphicMembership/FixedCoefficientFieldInComplex.lean}{493}{fixed frame values}
\releasecite{Development/AbelianVarietyTheory/Definitions/AutomorphicMembership/FixedCoefficientFieldInComplex.lean}{585}{the separate character}
\releasecite{Development/AbelianVarietyTheory/Proofs/IwasawaDecomposition/AddLargeBaseCoreSlot.lean}{22}{the preserved fields}.

The list field is contained in the compositum of \(M'\), the ray
field, the cyclotomic field of order \(N\) and the power-ideal Kummer
field \(V_{h_*}\). Here \(N\) is the least common multiple of
\(3^4 5^4\) and \(\prod_{a=0}^1|\mathrm{GL}_r(\mathbb F_{g_a})|^f\),
and \(h_*\) is the relative ray degree. The Kummer factor contains
the prescribed ideal and type-norm generators.
The even exponents of \(\mu_{0,j}\), the half shift of \(\eta\) and the
support and tameness of \(\delta,\delta'\) put the required readings of
these characters in that factor. The \(\omega\)-values are cyclotomic.
Put \(L=\log b=x^\varepsilon\) and \(r=k+m-1\).
The host bounds give \(\log(Qp_{\rm de}p'_{\rm de})\le20L+2\).
The ray exact sequence and \(\log h_E\le10 fL\) then give
\(\log h_*\le51 fL\). The finite level orders give
\(\log N\le19fr^2L+24\). The Kummer bound adds at most
\((u_E+2)\log h_*+\log h_E\) to the corresponding bound for \(E\),
where \(u_E\le2f\) is the unit rank. The CM bound uses the marked
quadratic discriminant parameter with logarithm at most \(4L\).
Adding the four degree and discriminant budgets and using \(f\ge8\)
and \(L\ge4\) bounds both logarithms by
\[
 19fr^2L+102f^2L+163fL+36L+46\le P_O=19fr^2L+198f^2L.
\]
Degree and root discriminant decrease on taking a subfield, so
\(\max\{\log[K_s:\mathbb Q],\log\operatorname{rd}(K_s)\}\le P_O\).
The point's field comparison applies this bound to the embedded
coefficient field of both runs and hence to the linked row summand.
\releasecite{Development/AbelianVarietyTheory/Proofs/IwasawaDecomposition/ClassNumberDvdHstar.lean}{347}{Fixed character containment}
\releasecite{Development/AbelianVarietyTheory/Proofs/IwasawaDecomposition/DegreeRootDiscrBoundKummerFieldOfPowerIdealRoots.lean}{25}{power-ideal Kummer bound}
\releasecite{Development/AbelianVarietyTheory/Proofs/IwasawaDecomposition/OmegaFieldDegreeRootDiscrBoundedByRingPolynomial.lean}{52}{four size budgets}
\releasecite{Development/AutomorphicArithmetic/ThetaCorrespondence/ThetaPairLifts/ThirdPlaceSeries/RowEPartTies.lean}{22}{field bound on the linked summand}.

At stage \(j\) the geometric cusp space \(C_{j,w}\) consists of the
framed section rules of weight \(w\) on this core over \(K_s\).
At a base point, read a short rule on the lift to its own marked
cover. The finite character and harmonic transformation law give the
same corrected value under a change of lift and the required covariance
under a change of frame. Uniqueness makes this reading natural under
base change and gives its descent to a section over \(K_s\). On the
marked models with the prescribed short-cover laws, chart and
component compatibilities, the positive last entry and copy number
at least \(8fr+4\) give complex spanning at the three short weights.
The readings of a \(K_s\)-basis are
linearly independent over \(\mathbb C\), so a proper \(K_s\)-span of
the short rules would remain proper after taking complex readings.
The finite lists in the
\hyperlink{proof.run-short-section-generators}{generator estimates}
therefore belong to this model. Their scaled norms and the integral
full list at \(w_L\) use the same field bound at either cover.
\releasecite{Formalism/LocalToGlobal/ProjectiveLattices/RecPJThAllNeBot.lean}{28}{Full and short rules of the actual model}
\releasecite{Development/AbelianVarietyTheory/Proofs/CuspFormEvaluations/ShortValTotOfPinned.lean}{62}{descent of the short reading}
\releasecite{Development/AbelianVarietyTheory/Proofs/SiegelThetaTransformation/ShortThetaSeriesSpanCuspFormsWeight.lean}{17}{stable-range spanning}
\releasecite{Development/AbelianVarietyTheory/Proofs/CuspFormEvaluations/RunScalarsNonzeroBoundedSupBound.lean}{52}{bounded integral lists}.

Write \(D_{\rm fl}=f\binom r2\), \(T_{\rm sec}=\lceil x^{377}\rceil\),
\(w_0=w_F+D_{\rm fl}w_L\), \(\kappa_j=w_F+T_{\rm sec} w_L\) and
\(C_j=C_{j,\kappa_j}\). On the run with own prime \(p_o\), let
\(\mathcal T_{j,p_o}\) be the inverse image under the weight inclusion
of \(C_j\) into the section algebra of the three scaled products
\[
 (A'_L)^{T_{\rm sec}}S'_F
      +(A'_L)^{T_{\rm sec}-D_{\rm fl}}S'_0+S'_{\kappa_j}.
\]
The \hyperlink{proof.geometric-section-index}{section index theorem}
puts this inverse image inside the canonical lattice \(M_{0,j,\kappa_j}\)
and a scalar multiple of the unscaled short lattice \(S_{\kappa_j,p_o}\).
Its hypotheses use the short-cover laws, the weight and growth
estimates and the prescribed chart data of this core.
\releasecite{Development/AbelianVarietyTheory/Proofs/CuspFormEvaluations/RunHypothesesFormsScalarsOfParameterConditions.lean}{49}{Scaled modules at the common weights}
\releasecite{Development/AbelianVarietyTheory/Proofs/DivisorProofs/RunsHoldOfTies.lean}{20}{both applications of the section theorem}
\releasecite{Development/AutomorphicArithmetic/ThetaCorrespondence/ThetaPairLifts/ThirdPlaceSeries/RunModelArch.lean}{37}{the actual field bound on the model}.

Put \(n_j=\dim_{K_s} C_j\). The marked
comparison matches the finite test points and their individual norms
on \(C_j\). Their unit balls intersect in the canonical lattice. The
two refinements use one isomorphism at this stage,
\(z_j:K^{n_j}\xrightarrow{\sim}K\otimes_{K_s}C_j\), and the same
base level action. For an integral module \(B\subset C_j\), write
\(\operatorname{pl}_{z_j}(B)=\operatorname{span}_{\mathcal O_K}
\{z_j^{-1}(1\otimes b):b\in B\}\). In the row coordinates the
comparison gives
\[
 \begin{aligned}
 M_j &= \operatorname{pl}_{z_j}(M_{0,j,\kappa_j}), &
 L_{vj} &= \operatorname{pl}_{z_j}(\mathcal T_{j,p_o(v)}),\\
 L_{K,vj} &= \operatorname{pl}_{z_j}(S_{\kappa_j,p_o(v)}), &
 \sigma_{vj} &= \text{the prescribed section scalar in }K^\times.
 \end{aligned}
\]
Here \(L_{K,vj}\) is the short lattice and \(v\) selects the run.
Localizing gives the comparison over \(\mathcal O_v\). Write \(e_j\)
for cusp reading into the \(K_s\)-form of the row space. The period
coordinate map is \((1\otimes e_j)\circ z_j\), using the identification
of that form after extending scalars to \(K\). It intertwines the level
action and good projector, identifies the good character cut of \(M_j\)
with \(J_j^{\rm row}\), and gives the base and twisted theta-column
lattices in the same coordinates. The rational class-sum identity
uses its original class-sum hypothesis.
\releasecite{Formalism/LocalToGlobal/IndefiniteReal/FWTie.lean}{17}{Equality of finite tests and norms}
\releasecite{Development/AbelianVarietyTheory/Proofs/DivisorProofs/NotOntoOfCharAlone.lean}{550}{common comparison for both runs}
\releasecite{Development/AbelianVarietyTheory/Definitions/Jets/EmbShortThLeThetaPlain.lean}{57}{placed lattices, action and scalar}
\releasecite{Development/AbelianVarietyTheory/Definitions/Foundations/ThetaSeriesTies.lean}{155}{period coordinates and class sums}
\releasecite{Development/AbelianVarietyTheory/Proofs/DivisorProofs/NotOntoOfCharAlone.lean}{445}{good row lattice}.

At a finite place \(a\) of \(K_s\) on either run, let \(E_a\) be
the prescribed ceiling, which bounds the product depth in the
geometric index theorem. In a local Smith basis write \(b_l\) for
the elementary exponents of the product lattice in the canonical
lattice. Apply the all-subspace bound at \(a\) to the span of the
basis vectors with \(b_l>E_a\). The bound gives
\(\sum_l(b_l-E_a)_+\le r_{aj}\), where \(r_{aj}\ge0\) is the
geometric remainder for this run and stage. On extending to \(v\mid a\)
both the exponents and the ceiling multiply by \(e(v/a)\). Write
\(E_{vj}=e(v/a)E_a\). The lattice
\(M_{j,v}\cap\varpi_v^{-E_{vj}}L_{vj,v}\) consequently has index
length at most \(e(v/a)r_{aj}\) in \(M_{j,v}\). For any \(K\)-subspace
\(V'\), work in its local extension \(V'_v\). The quotient of
\(M_{j,v}\cap V'_v\) by its intersection with this lattice injects into
the same finite quotient. Multiplying that intersection by
\(\varpi_v^{E_{vj}}\) puts it in \(L_{vj,v}\cap V'_v\) at additional
cost \(E_{vj}\dim_K V'\). Thus the same remainders give
\[
 \begin{gathered}
 \operatorname{length}_{\mathcal O_v}
   \big((M_{j,v}\cap V'_v)/(L_{vj,v}\cap V'_v)\big)
     \le E_{vj}\dim_K V'+\rho_{vj}^{(p_o(v))},\\
 \sum_{v\in A}\lambda_v\rho_{vj}^{(o)}\le n_j fmb\,d_2
    \quad (A\text{ finite},\ o=p_{\rm de},p'_{\rm de}).
 \end{gathered}
\]
The two budgets follow from the geometric sums and the local degree
formula \(\sum_{v\mid a}\lambda_v e(v/a)=\lambda_a\) with
\(\lambda_v=\log(\mathrm{N}v)/[K:\mathbb Q]\). They cover all cuts,
including the character cuts defined over \(K\). In the signed index
for the short lattice, \(L_{vj}\subseteq\sigma_{vj}L_{K,vj}\)
subtracts \(\operatorname{ord}_v(\sigma_{vj})\dim_K V'\). The depth
comparison uses \((E_{vj}-\operatorname{ord}_v\sigma_{vj})/e_v\) on
that run, with \(e_v\) the absolute ramification index. The two own-prime sums
and the conjugate pairing contribute at most
\(4n_j fmb/(\delta d)\le12 fmb\) per stage to the normalized
determinant estimate, since \(d_2\le1\) and \(n_j/\delta<3d\).
Here \(d\) is the good character-cut dimension and \(\delta\) the
index of the upper unipotent subgroup in the finite mirabolic group.
\releasecite{Development/AbelianVarietyTheory/Proofs/DivisorProofs/HasIndexBoundOfRuns.lean}{95}{All cuts after extension}
\releasecite{Development/AbelianVarietyTheory/Proofs/DivisorProofs/HasIndexBoundOfRuns.lean}{129}{the two separate budgets}
\releasecite{Development/AutomorphicArithmetic/ThetaCorrespondence/ThetaPairLifts/ThirdPlaceSeries/GeoStOfPageSt.lean}{29}{index and depth on the selected point}.
See the \hyperlink{proof.near-level-transport}{determinant calculation}.

At every embedding \(\iota:K\to\mathbb C\), the corresponding
framed readings give the supremum and Petersson norms on the
geometric model. The row comparison uses the extremal evaluation
when the restriction of \(\iota\) is good and the conjugate evaluation
when its conjugate is good. The common-measure identity gives, on
every cusp vector \(F\) in the selected model with the row norm taken
in its chosen embedding,
\[
 \|e_j F\|_{{\rm row},\iota}
    =\sqrt{\operatorname{rat}_j}\,\|F\|_{{\rm Pet},\iota|_{K_s}}.
\]
The ratio \(\operatorname{rat}_j>0\) is evaluated on the chosen theta
data. Its two logarithmic bounds give the required absolute bound on
the logarithm of the norm factor at the common threshold. By polarization
the same squared-norm comparison identifies the Gram forms. On the
primitive good space the pullback of \(J_{\rm good}\) to the stage
frame preserves its finite components and scales its vector norm by
\(n_{\rm lev}^{-1/2}\), where \(n_{\rm lev}\) is the order of the
base-face level group. Let \(M_j^{\rm main}\) denote the stage main
term of \eqref{eq:adjoint-stage-covolume}. In squared-slope units the
normalized log covolume satisfies
\[
 \frac{\operatorname{lcv}(J_j^{\rm row})}{d_{\rm row}}
     =-\mu(J_{\rm good})-\log n_{\rm lev}
     \le M_j^{\rm main}+C_3E_{\rm cov} .
\]
The last inequality uses the retained ordinary and terminal estimates
for this good lattice. It gives the \(L_{3,j}\) term in the mean.
The principal-level \(P_0,P_1\) character cuts then use their metric factors in
\eqref{eq:actual-row-level-metric} and the common squared-metric
compensation of the selected bundle.
\releasecite{Formalism/LocalToGlobal/ProjectiveLattices/XcLat.lean}{40}{Geometric norms at all embeddings}
\releasecite{Development/AutomorphicArithmetic/ThetaCorrespondence/ThetaPairLifts/ThirdPlaceSeries/ExactAtOfRows.lean}{64}{both evaluation signs}
\releasecite{Development/AutomorphicArithmetic/ThetaCorrespondence/ThetaPairLifts/ThirdPlaceSeries/CoveringMeasureReductionBound.lean}{18}{positive measure ratio and its bounds}
\releasecite{Development/AutomorphicArithmetic/ThetaCorrespondence/ThetaPairLifts/ThirdPlaceSeries/RecAtOfGramAt.lean}{17}{the frame and scaling}
\releasecite{Development/AutomorphicArithmetic/ThetaCorrespondence/ThetaPairLifts/ThirdPlaceSeries/EventualL3OfCovolume.lean}{14}{the covolume on the selected summands}.
See the \hyperlink{proof.absolute-stage-row-covolume}{stage covolume}
and \hyperlink{proof.row-level-mean}{compensated mean}.

\rhrange[Stronger compatibility from the pinned construction]{The
selected common core realizes both section runs over one small
embedded field at the index in question.}
{The cut bounds and both weighted budgets for the placed lattices,
the short-lattice inclusion with its scalar and depth, and the
covolume bounds in the normalized stage frames. This derivation
uses the common core, including the level action, projected tests
and period coordinates.}
The trace also uses the exact cloud and letters to obtain its
\hyperlink{proof.trace-word-ratios}{Vandermonde ratios}. The good
character share and discarded-type estimates use the marked two-level
action. At an ordinary stage the shifted rank comparison uses the
archimedean two-level calculation. At the terminal stage the free
level action gives the regular representation. Both comparisons
respect the character equations and good projector. The surjectivity
argument combines the index, depth, rank, share, twist and period
conditions on the same runs and coordinate object. Its norm estimate
enters the row proposition alongside the separate lower-eigenvalue
scalar estimate, the covolume, the functional and remaining row
bounds, and the definite-side condition. The original representation,
trace, period and numerical inputs continue to apply.
\releasecite{Development/AutomorphicArithmetic/ThetaCorrespondence/ThetaPairLifts/ThirdPlaceSeries/RunModelArch.lean}{87}{Archimedean two-level comparison}
\releasecite{Development/AbelianVarietyTheory/Proofs/DivisorProofs/BadWStOfRegular.lean}{40}{regular and shifted rank comparisons}
\releasecite{Development/AutomorphicArithmetic/ThetaCorrespondence/ThetaPairLifts/ThirdPlaceSeries/OntoPath.lean}{49}{the joint application}
\releasecite{Development/AutomorphicArithmetic/ThetaCorrespondence/ThetaPairLifts/ThirdPlaceSeries/LinkROfIsLinked.lean}{80}{row proposition carried by the selector}.

This choice supplies the section systems for the simultaneous
existence argument. A separately specified geometric model requires
its own comparison of canonical and short lattices, marked evaluations,
level actions and metrics, or bounds for their discrepancies within
the same finite and archimedean allowances. Even rank-one lattices
\(\mathcal O_v\) and \(\varpi_v^a\mathcal O_v\) over the same field
and in the same rational span differ by arbitrary local length \(a\).
